O conteúdo superior de um subconjunto \(E\) do plano \(\mathbb{R}^{2}\) é definido como
\[\mathcal{C}(E)=\inf\left\{\sum_{i=1}^{n} \operatorname{diâmetro}(E_{i})\right\}\] onde \(\inf\) é tomado sobre todas as famílias finitas de conjuntos \(E_{1},\dots,E_{n}\), \(n\in \mathbb{N}\), em \(\mathbb{R}^{2}\) tais que \(E\subset \bigcup_{i=1}^{n}E_{i}\).
O conteúdo inferior de \(E\) é definido como
\[\mathcal{K}(E)=\sup\left\{\operatorname{comprimento}(L) : \begin{array}{l} \text{\(L\) é um segmento de reta fechado} \\ \text{sobre o qual \(E\) pode ser contraído} \end{array} \right\}.\]
Prove que
\begin{itemize}
\item[(a)] \(\mathcal{C}(L)=\operatorname{comprimento}(L)\) se \(L\) é um segmento de reta fechado;
\item[(b)] \(\mathcal{C}(E) \geq \mathcal{K}(E)\);
\item[(c)] a igualdade em (b) não é sempre verdadeira mesmo se \(E\) for compacto.
\end{itemize}
